\section{Del perceptrón a las redes neuronales recurrentes}

\subsection{Repaso del perceptrón}

En la clase anterior, aprendimos la estructura básica del perceptrón: dados los insumos $x^{(1)}, x^{(2)}, \ldots, x^{(m)}$, cada insumo se multiplica por el peso correspondiente $w_1, w_2, \ldots, w_m$, se suman y el resultado pasa por una función de activación no lineal para producir la salida.

\begin{figure}[H]
\centering
\includegraphics[width=0.7\textwidth]{slides-images/slide-10.jpg}
\caption{Repaso de la estructura del perceptrón\protect\footnotemark}
\end{figure}
\footnotetext{Página 10 de las diapositivas.}

Expresado con fórmulas matemáticas:

$$z = \sum_{i=1}^{m} w_i x^{(i)} + b$$

$$\hat{y} = g(z)$$

Donde:
\begin{itemize}
    \item $x^{(i)}$: la $i$-ésima característica de entrada
    \item $w_i$: el peso correspondiente
    \item $b$: el término de sesgo
    \item $g(\cdot)$: la función de activación no lineal
    \item $\hat{y}$: la salida predicha
\end{itemize}

\subsection{Limitaciones del perceptrón para procesar datos secuenciales}

Cuando aplicamos el perceptrón (o una red feedforward) directamente a datos secuenciales, nos encontramos con un problema fundamental: \textbf{la predicción de cada paso de tiempo es completamente independiente}, sin aprovechar la información de los demás pasos de tiempo.

\begin{figure}[H]
\centering
\includegraphics[width=0.75\textwidth]{slides-images/slide-13.jpg}
\caption{Una red feedforward procesa cada paso de tiempo de forma independiente: no puede aprovechar la información histórica\protect\footnotemark}
\end{figure}
\footnotetext{Página 13 de las diapositivas.}

En concreto, si usamos una red feedforward para procesar cada paso de tiempo de forma independiente:

$$\hat{y}_t = f(x_t)$$

Entonces la predicción $\hat{y}_2$ del paso de tiempo $t=2$ depende únicamente del insumo $x_2$, ignorando por completo la información clave que podrían contener $x_0$ y $x_1$.

Esto es como observar solo la posición actual de una pelota sin tomar en cuenta su trayectoria anterior: la predicción es necesariamente ciega.

\subsection{La introducción de la recurrencia: la idea central del Recurrence}

Para resolver este problema, necesitamos un mecanismo que transmita al paso de tiempo actual la información de los pasos de tiempo pasados. La idea central es:

\begin{importantbox}{El núcleo de las redes neuronales recurrentes: el estado oculto}
Se define un \textbf{estado interno (Hidden State)} $h_t$, que se transmite entre los pasos de tiempo y conserva la ``memoria'' de la red sobre los insumos pasados. La predicción de cada paso de tiempo depende no solo del insumo actual $x_t$, sino también del estado oculto del paso anterior $h_{t-1}$:

$$\hat{y}_t = f(x_t, h_{t-1})$$
\end{importantbox}

\begin{figure}[H]
\centering
\includegraphics[width=0.85\textwidth]{slides-images/slide-15.jpg}
\caption{Una neurona recurrente: la salida depende del insumo actual y de la memoria pasada\protect\footnotemark}
\end{figure}
\footnotetext{Página 15 de las diapositivas.}

La imagen anterior muestra a la izquierda la vista plegada (Folded View) de la RNN y a la derecha la vista desplegada (Unrolled View). El estado oculto $h_t$ se transmite entre los pasos de tiempo como si fuera una cadena, lo que permite a la red ``recordar'' la información que vio previamente.

\subsection{Resumen de la sección}

La trayectoria del perceptrón a la RNN es clara: las redes feedforward carecen de la capacidad de procesar secuencias, por lo que introducimos el estado oculto como puente de información entre los pasos de tiempo, dando forma al concepto de recurrencia (Recurrence). Esta es la base para entender todos los modelos de secuencias.

\newpage
